\section{DA LEGITIMIDADE PARA REPRESENTAÇÃO}

\begin{enumerate}[resume]
\item A capacidade postulatória e os poderes de representação do subscritor estão devidamente comprovados pela procuração anexada (Anexo I), outorgada pelo Prefeito Municipal. Tal instrumento é plenamente suficiente para a atuação administrativa perante a distribuidora e a ANEEL, em conformidade com o art. 9º da REN nº 1.000/2021 e demais legislações aplicáveis.
\item Qualquer exigência documental que exceda a procuração, como a apresentação do contrato de honorários, não possui amparo legal e infringe as prerrogativas da advocacia, estabelecidas pela Lei nº 8.906/1994.
\item Ainda assim, com o intuito de evitar dilações indevidas e demonstrar total transparência, anexa-se voluntariamente cópia do contrato Nº \numeroContrato{} de prestação de serviços celebrado entre a municipalidade e a sociedade de advogados em \dataPublicacao{}.
\item O referido contrato, em sua \textbf{\tituloClausula{}}, confere poderes expressos e específicos para que este representante atue em nome do Município em todas as instâncias administrativas necessárias, notadamente perante a concessionária e a ANEEL, para a recuperação de valores e defesa de seus interesses.
\Citacao{\clausula{}}
\item A cláusula transcrita confere expressamente poderes de representação ao representante do escritório contratado pelo Município, nos moldes do Contrato de Mandato já juntado aos autos, o que torna indiscutível sua legitimidade para atuar em nome do ente municipal perante esta distribuidora.
\item Adicionalmente, cumpre informar que o contrato em questão foi devidamente publicado no(a) \textbf{\publicacao{}}, em obediência à legislação vigente, ato que lhe confere plena eficácia jurídica e oponibilidade contra terceiros.
\item Dessa forma, o requisito formal de publicação do referido instrumento contratual foi integralmente cumprido. 
\end{enumerate}
